% The General Sinus: Trigonometric Functions of an Oblique Projection
% Named version (rev3) of the Elemente der Mathematik manuscript.
\documentclass[11pt]{amsart}

\usepackage[a4paper,margin=3cm]{geometry}
\usepackage{amsmath, amssymb, mathtools}
\usepackage{tikz}
\usetikzlibrary{arrows.meta,calc,angles,quotes}
\usepackage{pgfplots}
\pgfplotsset{compat=1.17}
\usepackage{hyperref}

\newtheorem{thm}{Theorem}[section]
\newtheorem{cor}[thm]{Corollary}
\newtheorem{lem}[thm]{Lemma}
\newtheorem{prop}[thm]{Proposition}
\theoremstyle{definition}
\newtheorem{defn}[thm]{Definition}
\theoremstyle{remark}
\newtheorem{rem}[thm]{Remark}
\newtheorem*{ex}{Example}
\numberwithin{equation}{section}

\newcommand{\Com}{\mathrm{Com}}
\newcommand{\pb}{\overline{\phi}}

\begin{document}

\title{The General Sinus: Trigonometric Functions of an Oblique Projection}
\author{Mostafa Derraz}
\address{Independent Researcher, Tangier, Morocco}
\email{m.derraz@proton.me}

\subjclass[2020]{Primary 26A09; Secondary 51N20, 33B10}
\keywords{oblique projection, oblique coordinates, law of sines,
two-variable trigonometric functions, addition formulas, product identities}

\maketitle


% =======================================================================
\section{Introduction}
\label{sec:introduction}
% =======================================================================

The classical sine and cosine of an angle $\alpha$ come from one
construction. Take a point $Z$ on the second side of the angle,
project it straight down onto the first side, and read the sides of the
right triangle that results~\cite{Maor1998}. The right angle is part of
the construction, so familiar that it is rarely stated.

This paper changes one thing. The point $Z$ is projected onto the first
side along a line that makes a chosen angle $\phi$ with $ZO$, the
\emph{projection angle} (Figure~\ref{fig:projections}). The triangle
$OZ\overline{Z}$ is then oblique, with the angles $\alpha$ at $O$,
$\phi$ at $Z$ and $\pi-(\alpha+\phi)$ at the foot $\overline{Z}$. Its
sides, divided by $[OZ]$, define the functions $\sin(\alpha,\phi)$ and
$\cos(\alpha,\phi)$, and their quotient $\tan(\alpha,\phi)$. The law of
sines~\cite{CoxeterGreitzer1967} gives them at once,
\[
\sin(\alpha,\phi) = \frac{\sin\alpha}{\sin(\alpha+\phi)}, \qquad
\cos(\alpha,\phi) = \frac{\sin\phi}{\sin(\alpha+\phi)}, \qquad
\tan(\alpha,\phi) = \frac{\sin\alpha}{\sin\phi},
\]
and the classical functions return when $\alpha+\phi = \pi/2$.

These formulas are elementary, and we do not claim otherwise. What makes
them worth a systematic study is that they have a natural reading. When
$[OZ] = 1$, the projection line is parallel to a second axis through $O$
at the angle $\pb = \alpha+\phi$ to the first, and $\cos(\alpha,\phi)$
and $\sin(\alpha,\phi)$ are the coordinates of the point $Z$ of the unit
circle in the oblique frame of the two axes
(Proposition~\ref{prop:coordinates}). The classical functions are the
coordinates of the unit circle in an orthonormal frame; the general
functions are its coordinates in an oblique one. Oblique coordinates are
as old as analytic geometry, but we are not aware of a study of the
coordinate functions of the circle in such a frame as functions of two
angles in their own right. The closest work concerns the generalized
complex numbers, whose unit satisfies $i^2 = p + qi$
\cite{HarkinHarkin2004,KantorSolodovnikov1989,Yaglom1968}. In the
elliptic case these algebras are the plane named in an oblique frame and
the general functions are the coordinates of their unit circle; the
arithmetic of numbers written in this form is developed
in~\cite{DerrazGN}. The name \emph{generalized trigonometric
functions} is also used for the eigenfunctions of the one-dimensional
$p$-Laplacian~\cite{Lindqvist1995}; those functions are unrelated to
ours.

The general sine was introduced in~\cite{GeneralSinus}. The present
paper gives a self-contained account, organised around the coordinate
reading. Its main results are the following.
\begin{itemize}
\item In an oblique frame, the Pythagorean identity acquires one cross
  term, $\cos(\alpha,\phi)^2 + \sin(\alpha,\phi)^2
  + 2\cos(\alpha,\phi)\sin(\alpha,\phi)\cos(\alpha+\phi) = 1$
  (Proposition~\ref{prop:pythagoras}).
\item The addition formulas come in two kinds: with the projection angle
  fixed, so that the frame turns with the studied angle
  (Section~\ref{sec:addition}), and with the frame fixed
  (Proposition~\ref{prop:fixed-frame}). All of the first kind reduce to a
  single identity of Ptolemy; the second kind follows from Euler's
  formula written in the oblique frame and includes De Moivre's formula.
\item Differentiation with the frame held fixed acts on the general sine
  and cosine as a quarter turn,
  $D^n f(x,y) = f(x + n\pi/2,\, y - n\pi/2)$, so that $D^2 f = -f$ as for
  the classical functions (Theorem~\ref{thm:shift}).
\item Each general function is a product of $N$ copies of itself at scaled
  arguments, for every $N \ge 1$ (Theorem~\ref{thm:products}). Truncating
  each factor to second order gives rational functions which are again
  quotients of sines, in closed form, and which converge with an error of
  exact order $N^{-2}$ (Theorem~\ref{thm:GST-closed} and
  Corollary~\ref{cor:GST-rate}).
\end{itemize}
Section~\ref{sec:def} defines the functions and their domains,
Section~\ref{sec:coordinates} gives the coordinate reading, and
Section~\ref{sec:identities} collects symmetries and special values. The
results above occupy Sections~\ref{sec:addition}
to~\ref{sec:products}. Every statement is checked symbolically or at
$50$ significant digits by an accompanying script
(Section~\ref{sec:verification}).

\medskip
\noindent\textbf{Notation.} The studied angle is $\alpha$, or $x$ in
analytic statements; the projection angle is $\phi$, or $y$. Segment
lengths are written $[PQ]$. The classical functions keep their one-variable
notation, so that $\sin\alpha$ and $\sin(\alpha,\phi)$ are never confused.

% =======================================================================
\section{The general functions}
\label{sec:def}
% =======================================================================

\begin{figure}[t]
\centering
\begin{tikzpicture}[>=Stealth, scale=1.0, every node/.style={font=\small}]
  % (a) vertical projection
  \begin{scope}
    \coordinate (O) at (0,0);
    \coordinate (X) at (3.6,0);
    \coordinate (Z) at (35:3.0);
    \coordinate (Zb) at ($(O)!(Z)!(X)$);
    \draw[gray, ->] (-0.2,0) -- (3.8,0);
    \draw[thick] (O) -- (Z) -- (Zb) -- cycle;
    \draw ($(Zb)+(-0.22,0)$) -- ++(0,0.22) -- ++(0.22,0);
    \pic [draw, "$\alpha$", angle radius=0.75cm, angle eccentricity=1.3] {angle = X--O--Z};
    \fill (Z) circle (1.3pt) node[above] {$Z$};
    \fill (Zb) circle (1.3pt) node[below] {$\overline{Z}$};
    \node[below left] at (O) {$O$};
    \node at (1.8,-0.95) {(a) straight down};
  \end{scope}
  % (b) projection at angle phi
  \begin{scope}[xshift=5.4cm]
    \coordinate (O) at (0,0);
    \coordinate (X) at (3.6,0);
    \coordinate (Z) at (35:3.0);
    \coordinate (Zb) at ($(Z)+(255:{3.0*sin(35)/sin(75)})$);
    \draw[gray, ->] (-0.2,0) -- (3.8,0);
    \draw[thick] (O) -- (Z) -- (Zb) -- cycle;
    \pic [draw, "$\alpha$", angle radius=0.75cm, angle eccentricity=1.3] {angle = X--O--Z};
    \pic [draw, "$\phi$", angle radius=0.55cm, angle eccentricity=1.45] {angle = O--Z--Zb};
    \fill (Z) circle (1.3pt) node[above] {$Z$};
    \fill (Zb) circle (1.3pt) node[below] {$\overline{Z}$};
    \node[below left] at (O) {$O$};
    \node at (1.8,-0.95) {(b) at the angle $\phi$};
  \end{scope}
\end{tikzpicture}
\caption{Two projections of a point $Z$ onto the first side of the angle
$\alpha$. (a) Straight down: a right triangle, whose sides give $\cos\alpha$
and $\sin\alpha$. (b) Along a line making the angle $\phi$ with $ZO$: an
oblique triangle with the angles $\alpha$, $\phi$ and
$\pi-(\alpha+\phi)$, whose sides give $\cos(\alpha,\phi)$ and
$\sin(\alpha,\phi)$. Here $\alpha = 35^\circ$ and $\phi = 40^\circ$.}
\label{fig:projections}
\end{figure}

Let $\alpha$ be an angle with vertex $O$, let $Z \neq O$ be a point on
its second side, and let $0 < \phi < \pi - \alpha$. The line through $Z$
that makes the angle $\phi$ with $ZO$ meets the first side at a point
$\overline{Z}$, the \emph{projection of $Z$ at the angle $\phi$}. The
triangle $OZ\overline{Z}$ has the angles $\alpha$, $\phi$ and
$\pi-(\alpha+\phi)$. All the triangles obtained in this way from points
$Z$ on the same side are similar, so the ratios of their sides depend on
$\alpha$ and $\phi$ only.

\begin{defn}[General trigonometric functions]
\label{defn:general}
The \emph{general sine, cosine and tangent} are the ratios
\[
\sin(\alpha,\phi) = \frac{[Z\overline{Z}]}{[OZ]}, \qquad
\cos(\alpha,\phi) = \frac{[O\overline{Z}]}{[OZ]}, \qquad
\tan(\alpha,\phi) = \frac{[Z\overline{Z}]}{[O\overline{Z}]},
\]
and $\csc$, $\sec$, $\cot$ are their reciprocals.
\end{defn}

\begin{thm}[Analytic form]
\label{thm:analytic}
For $0 < \alpha$, $0 < \phi$ and $\alpha + \phi < \pi$,
\begin{equation}
\label{eq:analytic}
\sin(\alpha,\phi) = \frac{\sin\alpha}{\sin(\alpha+\phi)}, \qquad
\cos(\alpha,\phi) = \frac{\sin\phi}{\sin(\alpha+\phi)}, \qquad
\tan(\alpha,\phi) = \frac{\sin\alpha}{\sin\phi}.
\end{equation}
\end{thm}

\begin{proof}
By the law of sines in $OZ\overline{Z}$,
$[Z\overline{Z}]/\sin\alpha = [O\overline{Z}]/\sin\phi
= [OZ]/\sin(\pi-\alpha-\phi)$, and $\sin(\pi-\alpha-\phi) = \sin(\alpha+\phi)$.
\end{proof}

From here on, \eqref{eq:analytic} \emph{defines} the general functions
for all real $\alpha$ and $\phi$ at which the denominators do not vanish:
$\sin(\alpha,\phi)$ and $\cos(\alpha,\phi)$ need $\alpha+\phi \notin \pi\mathbb{Z}$,
and $\tan(\alpha,\phi)$ needs $\phi \notin \pi\mathbb{Z}$. The reciprocals are
\[
\csc(\alpha,\phi) = \frac{\sin(\alpha+\phi)}{\sin\alpha}, \qquad
\sec(\alpha,\phi) = \frac{\sin(\alpha+\phi)}{\sin\phi}, \qquad
\cot(\alpha,\phi) = \frac{\sin\phi}{\sin\alpha},
\]
defined where their denominators do not vanish (for $\csc$ and $\sec$ we
also keep $\alpha+\phi \notin \pi\mathbb{Z}$, so that they are reciprocals
of finite non-zero values).
For $\alpha$ and $\phi$ outside the triangle range, the general functions
keep the geometric meaning of Proposition~\ref{prop:coordinates}, as signed
coordinates.

\begin{cor}[Classical case]
\label{cor:classical}
If $\alpha + \phi = \pi/2$, then
\[
\sin(\alpha,\phi) = \sin\alpha, \qquad \cos(\alpha,\phi) = \cos\alpha, \qquad
\tan(\alpha,\phi) = \tan\alpha .
\]
\end{cor}

\begin{proof}
Then $\sin(\alpha+\phi) = 1$ and $\sin\phi = \cos\alpha$.
\end{proof}

The condition $\alpha+\phi = \pi/2$ says that the triangle has its right
angle at $\overline{Z}$. The general functions therefore form a
one-parameter family through the classical ones, the parameter recording
the direction of projection.

\begin{ex}
For $\alpha = \pi/4$ and $\phi = \pi/3$,
\begin{gather*}
\sin\Bigl(\frac{\pi}{4},\frac{\pi}{3}\Bigr)
= \frac{\sin(\pi/4)}{\sin(7\pi/12)}
= \frac{\sqrt2/2}{(\sqrt6+\sqrt2)/4} = \sqrt3 - 1 \approx 0.732, \\
\cos\Bigl(\frac{\pi}{4},\frac{\pi}{3}\Bigr)
= \frac{\sqrt3/2}{(\sqrt6+\sqrt2)/4} = \frac{3\sqrt2-\sqrt6}{2} \approx 0.897.
\end{gather*}
\end{ex}

\begin{figure}[t]
\centering
\begin{tikzpicture}
\begin{axis}[
  width=0.92\textwidth, height=6.2cm,
  xmin=0, xmax=180, ymin=-3, ymax=3,
  xtick={0,30,60,90,120,150,180},
  xlabel={$\alpha$ (degrees)}, ylabel={$\sin(\alpha,\phi)$},
  axis lines=middle, every axis x label/.style={at={(current axis.south east)},anchor=north east, yshift=-2pt},
  every axis y label/.style={at={(current axis.north west)}, anchor=south west},
  legend style={at={(0.5,-0.14)}, anchor=north, font=\footnotesize, draw=none, legend columns=2, column sep=8pt},
  unbounded coords=jump, restrict y to domain=-3.2:3.2,
  samples=400, domain=0.5:179.5,
  tick label style={font=\footnotesize}, label style={font=\small},
]
  \addplot[thick, black, densely dotted] {sin(x)/sin(x+90)};
  \addlegendentry{$\phi = 90^\circ$}
  \addplot[thick, black] {sin(x)/sin(x+60)};
  \addlegendentry{$\phi = 60^\circ$}
  \addplot[thick, gray] {sin(x)/sin(x+30)};
  \addlegendentry{$\phi = 30^\circ$}
  \addplot[thick, black, dashed, domain=0.5:179.5] {sin(x)/sin(x+90-x)};
  \addlegendentry{$\phi = 90^\circ - \alpha$ (classical $\sin\alpha$)}
\end{axis}
\end{tikzpicture}
\caption{The general sine as a function of $\alpha$ for three fixed
projection angles. It has a pole at $\alpha = 180^\circ - \phi$, where
the projection line becomes parallel to the first side. Along the line
$\alpha+\phi = 90^\circ$ (dashed) it is the classical sine. For
$\phi = 90^\circ$ it is $\tan\alpha$ (dotted), by
Proposition~\ref{prop:special}.}
\label{fig:plot}
\end{figure}

% =======================================================================
\section{The coordinate reading}
\label{sec:coordinates}
% =======================================================================

Draw the unit circle about $O$ together with a second axis through $O$,
the \emph{composer axis} $\Com$, making the angle $\pb$ with the
$x$-axis, $0 < \pb < \pi$ (Figure~\ref{fig:composed}). We call this
configuration the \emph{composed circle}. For the point $Z$ of the circle
at the angle $\alpha$, projecting $Z$ onto the $x$-axis parallel to the
composer axis is the projection at the angle $\phi = \pb - \alpha$, since
the projection line and $ZO$ meet at that angle. The three angles are
therefore tied by
\begin{equation}
\label{eq:phibar}
\pb = \alpha + \phi .
\end{equation}

\begin{figure}[t]
\centering
\begin{tikzpicture}[>=Stealth, scale=2.4, every node/.style={font=\small}]
  \def\al{30}
  \def\pbar{70}
  \coordinate (O) at (0,0);
  \coordinate (X) at (1.35,0);
  \coordinate (Z) at (\al:1);
  \coordinate (Zb) at ($(Z)+(\pbar+180:{sin(\al)/sin(\pbar)})$);
  \coordinate (Com) at (\pbar:1.32);
  \draw[gray!60] (O) circle (1);
  \draw[->, gray] (-1.25,0) -- (1.45,0) node[right] {$x$};
  \draw[->, thick] (O) -- (Com) node[above] {$\Com$};
  \draw[thick] (O) -- (Z);
  \draw[very thick] (Zb) -- (Z);
  \draw[very thick, gray] (O) -- (Zb);
  \pic [draw, "$\alpha$", angle radius=0.42cm, angle eccentricity=1.35] {angle = X--O--Z};
  \draw[gray] (0.52,0) arc (0:\pbar:0.52);
  \node[gray] at (56:0.64) {$\overline{\phi}$};
  \pic [draw, "$\phi$", angle radius=0.28cm, angle eccentricity=1.75] {angle = O--Z--Zb};
  \fill (Z) circle (0.7pt) node[above right] {$Z$};
  \fill (Zb) circle (0.7pt) node[below] {$\overline{Z}$};
  \node[below left] at (O) {$O$};
  \node[below] at ($(O)!0.5!(Zb)+(0,-0.08)$) {$\cos(\alpha,\phi)$};
  \node[right] at ($(Zb)!0.55!(Z)+(0.03,0)$) {$\sin(\alpha,\phi)$};
\end{tikzpicture}
\caption{The composed circle. The point $Z$ of the unit circle at the
angle $\alpha$, projected onto the $x$-axis parallel to the composer axis
($\pb = \alpha+\phi$), has the coordinate $\cos(\alpha,\phi)$ along the
$x$-axis and $\sin(\alpha,\phi)$ along the composer axis
(Proposition~\ref{prop:coordinates}). Here $\alpha = 30^\circ$ and
$\pb = 70^\circ$.}
\label{fig:composed}
\end{figure}

\begin{prop}[Coordinates in an oblique frame]
\label{prop:coordinates}
Let $0 < \pb < \pi$, let $\vec u = (\cos\pb, \sin\pb)$ be the unit vector
of the composer axis, and let $Z = (\cos\alpha, \sin\alpha)$, with
$\phi = \pb - \alpha$. Then
\[
\overrightarrow{OZ} = \cos(\alpha,\phi)\,\vec e_1 + \sin(\alpha,\phi)\,\vec u ,
\]
where $\vec e_1 = (1,0)$. That is, $\cos(\alpha,\phi)$ and $\sin(\alpha,\phi)$
are the coordinates of $Z$ in the frame $(\vec e_1, \vec u)$, for every
real $\alpha$.
\end{prop}

\begin{proof}
Solve $(\cos\alpha,\sin\alpha) = a(1,0) + b(\cos\pb,\sin\pb)$. The second
components give $b = \sin\alpha/\sin\pb$, and the first give
$a = \cos\alpha - b\cos\pb = \sin(\pb-\alpha)/\sin\pb$. With
$\pb = \alpha+\phi$ these are $\sin(\alpha,\phi)$ and $\cos(\alpha,\phi)$.
Since $\sin\pb \neq 0$, the solution is unique.
\end{proof}

The proposition gives the general functions a meaning for all real
$\alpha$, as signed coordinates, and it explains their poles: as
$\pb \to 0$ or $\pi$ the composer axis falls onto the $x$-axis and the
oblique frame degenerates. The length of $\overrightarrow{OZ}$ written in
the frame gives the oblique form of Pythagoras.

\begin{prop}[Pythagorean identity of an oblique frame]
\label{prop:pythagoras}
Whenever $\alpha+\phi \notin \pi\mathbb{Z}$,
\begin{equation}
\label{eq:pythagoras}
\cos(\alpha,\phi)^2 + \sin(\alpha,\phi)^2
+ 2\cos(\alpha,\phi)\sin(\alpha,\phi)\cos(\alpha+\phi) = 1 .
\end{equation}
\end{prop}

\begin{proof}
$|a\vec e_1 + b\vec u|^2 = a^2 + b^2 + 2ab\cos\pb$ and $|\overrightarrow{OZ}| = 1$.
Equivalently, \eqref{eq:pythagoras} is the law of cosines at
$\overline{Z}$ in the triangle $OZ\overline{Z}$, whose angle there is
$\pi - \pb$.
\end{proof}

For $\pb = \pi/2$ the cross term vanishes and \eqref{eq:pythagoras} is
$\cos^2\alpha + \sin^2\alpha = 1$. In general the point
$(\cos(\alpha,\phi), \sin(\alpha,\phi))$, plotted with its two
coordinates at right angles, runs along an ellipse as $\alpha$ varies
with $\pb$ fixed; its axes lie on the diagonals~\cite{DerrazGN}.

% =======================================================================
\section{Symmetries and special values}
\label{sec:identities}
% =======================================================================

All statements in this section follow directly from \eqref{eq:analytic};
we indicate the one-line arguments.

\begin{prop}[Exchange and reciprocals]
\label{prop:exchange}
On the respective domains,
\[
\sin(\alpha,\phi) = \cos(\phi,\alpha), \qquad
\tan(\alpha,\phi) = \cot(\phi,\alpha), \qquad
\csc(\alpha,\phi) = \sec(\phi,\alpha),
\]
and the reciprocals are general functions themselves,
\begin{gather*}
\csc(\alpha,\phi) = \sin(\alpha+\phi,\,-\phi), \qquad
\sec(\alpha,\phi) = \tan(\alpha+\phi,\,\phi), \\
\cot(\alpha,\phi) = \cos(\alpha-\phi,\,\phi).
\end{gather*}
\end{prop}

\begin{proof}
Exchanging $\alpha$ and $\phi$ in \eqref{eq:analytic} exchanges the
numerators $\sin\alpha$ and $\sin\phi$. For the second line,
$\sin(\alpha+\phi,-\phi) = \sin(\alpha+\phi)/\sin\alpha$,
$\tan(\alpha+\phi,\phi) = \sin(\alpha+\phi)/\sin\phi$ and
$\cos(\alpha-\phi,\phi) = \sin\phi/\sin\alpha$.
\end{proof}

\begin{prop}[Signs and shifts]
\label{prop:symmetry}
On the respective domains, for all integers $p$ and $q$,
\begin{gather*}
\sin(-\alpha,-\phi) = \sin(\alpha,\phi), \qquad
\cos(-\alpha,-\phi) = \cos(\alpha,\phi), \\
\sin(-\alpha,\phi) = \sin(\alpha,-\phi), \\
\sin(\alpha+p\pi,\,\phi+q\pi) = (-1)^q\sin(\alpha,\phi), \\
\cos(\alpha+p\pi,\,\phi+q\pi) = (-1)^p\cos(\alpha,\phi), \\
\tan(\alpha+p\pi,\,\phi+q\pi) = (-1)^{p+q}\tan(\alpha,\phi).
\end{gather*}
\end{prop}

\begin{proof}
The sine is odd, and $\sin(t + m\pi) = (-1)^m \sin t$.
\end{proof}

In particular the general functions are $2\pi$-periodic in each variable.
The shift rule is not symmetric in the two variables: a half turn of the
projection angle changes the sign of the general sine, and a half turn of
the studied angle that of the general cosine.

\begin{prop}[Special values]
\label{prop:special}
On the respective domains,
\begin{gather*}
\sin(\alpha,\tfrac{\pi}{2}) = \tan\alpha, \qquad
\cos(\alpha,\tfrac{\pi}{2}) = \sec\alpha, \qquad
\tan(\alpha,\tfrac{\pi}{2}) = \sin\alpha, \\
\sin(\theta,\theta) = \cos(\theta,\theta) = \frac{1}{2\cos\theta}, \qquad
\tan(\theta,\theta) = 1 .
\end{gather*}
\end{prop}

\begin{proof}
$\sin(\alpha+\pi/2) = \cos\alpha$, and $\sin 2\theta = 2\sin\theta\cos\theta$.
\end{proof}

At $\theta = \pi/3$ all six general functions of $(\theta,\theta)$ equal
$1$: the triangle is equilateral. At $\theta = \pi/5$,
$\sin(\theta,\theta) = \cos(\theta,\theta) = 1/(2\cos 36^\circ) = \tau^{-1}$,
where $\tau = (1+\sqrt5)/2$ is the golden ratio, and the triangle is the
golden triangle with angles $36^\circ$, $36^\circ$, $108^\circ$.

% =======================================================================
\section{Addition formulas}
\label{sec:addition}
% =======================================================================

The addition formulas of the general functions all reduce, after
clearing denominators, to one identity among ordinary sines.

\begin{lem}[Ptolemy's identity]
\label{lem:ptolemy}
For all real $x$, $y$ and $z$,
\begin{equation}
\label{eq:ptolemy}
\sin(x+y)\sin(y+z) = \sin x\,\sin z + \sin y\,\sin(x+y+z).
\end{equation}
\end{lem}

\begin{proof}
By $2\sin P\sin Q = \cos(P-Q) - \cos(P+Q)$, both sides equal
$\tfrac12\bigl[\cos(x-z) - \cos(x+2y+z)\bigr]$.
\end{proof}

For a cyclic quadrilateral with consecutive arcs $2x$, $2y$, $2z$ and
$2\pi - 2(x+y+z)$, inscribed in a circle of diameter $1$, the chords are
$\sin x$, $\sin y$, $\sin z$, $\sin(x+y+z)$ and the diagonals
$\sin(x+y)$, $\sin(y+z)$, and \eqref{eq:ptolemy} is Ptolemy's theorem;
hence the name.

\begin{thm}[Addition in the studied angle]
\label{thm:add-studied}
On the domains where both sides are defined,
\begin{align}
\label{eq:add-sin}
\sin(\alpha+\beta,\phi) &= \sin(\alpha,\phi) + \cos(\alpha,\phi)\,\sin(\beta,\phi+\alpha), \\
\label{eq:add-cos}
\cos(\alpha+\beta,\phi) &= \cos(\alpha,\phi)\cos(\beta,\phi) + \sin(\alpha,\phi)\sin(\beta,\phi)\cos(\alpha+\beta,\phi), \\
\label{eq:add-tan}
\tan(\alpha+\beta,\phi) &= \tan(\alpha,\phi)\cos\beta + \tan(\beta,\phi)\cos\alpha .
\end{align}
\end{thm}

\begin{proof}
Multiply \eqref{eq:add-sin} by $\sin(\alpha+\phi)\sin(\alpha+\beta+\phi)$:
it becomes $\sin(\alpha+\beta)\sin(\alpha+\phi) = \sin\alpha\,\sin(\alpha+\beta+\phi)
+ \sin\phi\,\sin\beta$, which is \eqref{eq:ptolemy} with
$(x,y,z) = (\beta,\alpha,\phi)$. Multiply \eqref{eq:add-cos} by
$\sin(\alpha+\phi)\sin(\beta+\phi)\sin(\alpha+\beta+\phi)/\sin\phi$: it
becomes $\sin(\alpha+\phi)\sin(\beta+\phi) = \sin\phi\,\sin(\alpha+\beta+\phi)
+ \sin\alpha\,\sin\beta$, which is \eqref{eq:ptolemy} with
$(x,y,z) = (\alpha,\phi,\beta)$. Finally \eqref{eq:add-tan} is
$\sin(\alpha+\beta) = \sin\alpha\cos\beta + \cos\alpha\sin\beta$ divided
by $\sin\phi$.
\end{proof}

The cosine formula \eqref{eq:add-cos} contains $\cos(\alpha+\beta,\phi)$
on both sides. Solving for it, and inverting, gives two closed forms:
\begin{gather}
\label{eq:add-closed}
\cos(\alpha+\beta,\phi) = \frac{\cos(\alpha,\phi)\cos(\beta,\phi)}{1 - \sin(\alpha,\phi)\sin(\beta,\phi)},
\\
\label{eq:add-sec}
\sec(\alpha+\beta,\phi) = \sec(\alpha,\phi)\sec(\beta,\phi) - \tan(\alpha,\phi)\tan(\beta,\phi).
\end{gather}
The second is the exact analogue of $\cos(\alpha+\beta) = \cos\alpha\cos\beta - \sin\alpha\sin\beta$,
with the secant and tangent in the roles of the cosine and sine.

\begin{thm}[Addition in the projection angle]
\label{thm:add-projection}
On the domains where both sides are defined,
\begin{align}
\label{eq:addp-sin}
\sin(\alpha,\phi+\theta) &= \sin(\alpha,\phi)\,\sin(\alpha+\phi,\theta), \\
\label{eq:addp-cos}
\cos(\alpha,\phi+\theta) &= \cos(\alpha,\phi) + \sin(\alpha,\phi)\,\cos(\alpha+\phi,\theta).
\end{align}
\end{thm}

\begin{proof}
The right side of \eqref{eq:addp-sin} is
$\frac{\sin\alpha}{\sin(\alpha+\phi)}\cdot\frac{\sin(\alpha+\phi)}{\sin(\alpha+\phi+\theta)}$.
Multiplying \eqref{eq:addp-cos} by $\sin(\alpha+\phi)\sin(\alpha+\phi+\theta)$
gives $\sin(\alpha+\phi)\sin(\phi+\theta) = \sin\phi\,\sin(\alpha+\phi+\theta)
+ \sin\alpha\,\sin\theta$, which is \eqref{eq:ptolemy} with
$(x,y,z) = (\alpha,\phi,\theta)$.
\end{proof}

Formula \eqref{eq:addp-sin} composes two oblique projections: projecting
at the angle $\phi$ and then, from the new point, at the angle $\theta$,
multiplies the general sines, just as similar triangles multiply ratios.
Formulas \eqref{eq:add-sin}--\eqref{eq:add-tan} have a geometric proof
on the composed circle as well, by cutting the triangle of
$\alpha+\beta$ into the triangles of $\alpha$ and $\beta$; the algebraic
proof above is shorter and covers all signs at once.

\begin{cor}[Double and half angles]
\label{cor:double}
On the domains where both sides are defined,
\begin{gather*}
\sin(2\alpha,\phi) = \sin(\alpha,\phi)\bigl(1 + \cos(2\alpha,\phi)\bigr), \qquad
\cos(2\alpha,\phi) = \frac{\cos(\alpha,\phi)^2}{1 - \sin(\alpha,\phi)^2}, \\
\sec(2\alpha,\phi) = \sec(\alpha,\phi)^2 - \tan(\alpha,\phi)^2, \qquad
\tan(2\alpha,\phi) = 2\tan(\alpha,\phi)\cos\alpha ,
\end{gather*}
and, replacing $\alpha$ by $\alpha/2$,
$\cos(\alpha,\phi) = \cos(\tfrac{\alpha}{2},\phi)^2\big/\bigl(1 - \sin(\tfrac{\alpha}{2},\phi)^2\bigr)$.
\end{cor}

\begin{proof}
Put $\beta = \alpha$ in Theorem~\ref{thm:add-studied}, \eqref{eq:add-closed}
and \eqref{eq:add-sec}. For the first formula,
\[
\cos(\alpha,\phi)\sin(\alpha,\phi+\alpha)
= \frac{\sin\phi}{\sin(\alpha+\phi)}\cdot\frac{\sin\alpha}{\sin(2\alpha+\phi)}
= \sin(\alpha,\phi)\cos(2\alpha,\phi). \qedhere
\]
\end{proof}

% =======================================================================
\section{The fixed frame: multiplication and derivatives}
\label{sec:calculus}
% =======================================================================

In analytic statements we write $(x,y)$ for $(\alpha,\phi)$ and
$f_x$, $f_y$ for partial derivatives.

\begin{prop}[Partial derivatives]
\label{prop:partials}
On the domains of the functions,
\begin{align*}
\partial_x \sin(x,y) &= \frac{\cos(x,y)}{\sin(x+y)}, &
\partial_y \sin(x,y) &= -\sin(x,y)\cot(x+y), \\
\partial_x \cos(x,y) &= -\cos(x,y)\cot(x+y), &
\partial_y \cos(x,y) &= \frac{\sin(x,y)}{\sin(x+y)}, \\
\partial_x \tan(x,y) &= \frac{\cos x}{\sin y}, &
\partial_y \tan(x,y) &= -\tan(x,y)\cot y,
\end{align*}
and the second derivatives are
\begin{align*}
\partial_x^2 \sin(x,y) &= -2\cot(x+y)\,\partial_x\sin(x,y), \\
\partial_y^2 \sin(x,y) &= \sin(x,y)\bigl(1 + 2\cot^2(x+y)\bigr), \\
\partial_x^2 \cos(x,y) &= \cos(x,y)\bigl(1 + 2\cot^2(x+y)\bigr), \\
\partial_y^2 \cos(x,y) &= -2\cot(x+y)\,\partial_y\cos(x,y), \\
\partial_x^2 \tan(x,y) &= -\tan(x,y), \qquad
\partial_y^2 \tan(x,y) = \tan(x,y)\bigl(1 + 2\cot^2 y\bigr).
\end{align*}
\end{prop}

\begin{proof}
Direct differentiation of \eqref{eq:analytic}. For instance
$\partial_x\sin(x,y) = \bigl(\cos x\sin(x+y) - \sin x\cos(x+y)\bigr)/\sin^2(x+y)
= \sin y/\sin^2(x+y)$, and
$\partial_y^2 \bigl(\sin x/\sin(x+y)\bigr)
= \sin x\bigl(\sin^2(x+y) + 2\cos^2(x+y)\bigr)/\sin^3(x+y)$.
\end{proof}

The partial derivatives move the frame: changing $x$ with $y$ fixed
turns the composer axis, since $\pb = x+y$. So do the addition formulas of
Section~\ref{sec:addition}, in which the projection angle $\phi$ is fixed
and the frame angle $\pb = \alpha + \phi$ changes with $\alpha$. The rest of
this section keeps the frame fixed and moves the point along the circle.
For a fixed $\pb \in (0,\pi)$ write
\begin{equation}
\label{eq:CS}
C_\theta = \cos(\theta,\,\pb-\theta) = \frac{\sin(\pb-\theta)}{\sin\pb}, \qquad
S_\theta = \sin(\theta,\,\pb-\theta) = \frac{\sin\theta}{\sin\pb},
\end{equation}
the coordinates of the point of the unit circle at the angle $\theta$ in
the frame $(\vec e_1,\vec u)$ (Proposition~\ref{prop:coordinates}).
Identify the plane with $\mathbb{C}$, so that $\vec e_1 = 1$ and
$\vec u = u = e^{i\pb}$. Then Proposition~\ref{prop:coordinates} reads
\begin{equation}
\label{eq:oblique-euler}
e^{i\theta} = C_\theta + S_\theta\, u ,
\end{equation}
which is Euler's formula written in the oblique frame. Since $u$ is a root
of $t^2 - 2\cos\pb\,t + 1 = (t - e^{i\pb})(t - e^{-i\pb})$,
\begin{equation}
\label{eq:rule-square}
u^2 = 2\cos\pb\; u - 1 .
\end{equation}

\begin{prop}[Addition in a fixed frame]
\label{prop:fixed-frame}
For a fixed $\pb \in (0,\pi)$ and all real $\alpha$, $\beta$,
\begin{align}
\label{eq:ff-cos}
C_{\alpha+\beta} &= C_\alpha C_\beta - S_\alpha S_\beta , \\
\label{eq:ff-sin}
S_{\alpha+\beta} &= C_\alpha S_\beta + S_\alpha C_\beta + 2\cos\pb\; S_\alpha S_\beta ,
\end{align}
and for every integer $n$,
\begin{equation}
\label{eq:demoivre}
\bigl(C_\theta + S_\theta\,u\bigr)^n = C_{n\theta} + S_{n\theta}\,u .
\end{equation}
\end{prop}

\begin{proof}
Multiply $e^{i\alpha} = C_\alpha + S_\alpha u$ by $e^{i\beta} = C_\beta + S_\beta u$
and reduce $u^2$ with \eqref{eq:rule-square}:
\[
e^{i(\alpha+\beta)} = \bigl(C_\alpha C_\beta - S_\alpha S_\beta\bigr)
+ \bigl(C_\alpha S_\beta + S_\alpha C_\beta + 2\cos\pb\,S_\alpha S_\beta\bigr)\,u .
\]
Since $1$ and $u$ are linearly independent over $\mathbb{R}$, the coordinates
of $e^{i(\alpha+\beta)}$ in the frame are unique, and they are
$C_{\alpha+\beta}$ and $S_{\alpha+\beta}$ by \eqref{eq:oblique-euler}. Formula
\eqref{eq:demoivre} is $(e^{i\theta})^n = e^{in\theta}$ read in the same way.
\end{proof}

For $\pb = \pi/2$, \eqref{eq:ff-cos} and \eqref{eq:ff-sin} are the classical
addition formulas, and \eqref{eq:demoivre} is De Moivre's formula. In the
oblique frame the sine formula acquires the cross term
$2\cos\pb\,S_\alpha S_\beta$, the same cross term as in the Pythagorean
identity \eqref{eq:pythagoras}, which in this notation reads
$C_\theta^2 + 2\cos\pb\,C_\theta S_\theta + S_\theta^2 = 1$. With $\theta = \pb$
in \eqref{eq:demoivre}, $C_{n\pb} = -\sin((n-1)\pb)/\sin\pb$ and
$S_{n\pb} = \sin(n\pb)/\sin\pb = U_{n-1}(\cos\pb)$, a Chebyshev polynomial of
the second kind, so that the powers of $u$ in its own frame are
$u^n = -U_{n-2}(\cos\pb) + U_{n-1}(\cos\pb)\,u$ for $n \ge 2$.

The derivative that moves the point with the frame held fixed is the
derivative along the lines $x + y = \text{const}$.

\begin{defn}[Derivative in a fixed frame]
\label{defn:D}
For a differentiable function $f$ of $(x,y)$, let
\[
Df(x,y) = \frac{\mathrm{d}}{\mathrm{d}t}\, f(x+t,\,y-t)\Big|_{t=0}
= \partial_x f(x,y) - \partial_y f(x,y).
\]
\end{defn}

Along the path $t \mapsto (x+t, y-t)$ the studied angle grows while the
projection angle shrinks by the same amount, so $\pb = x+y$ is constant:
the point $Z$ turns on the unit circle and the oblique frame does not
move. This is the operation that the classical derivative of $\sin\alpha$
performs on the unit circle with its orthonormal frame.

\begin{thm}[Quarter-turn shift]
\label{thm:shift}
For every integer $n \ge 0$ and $(x,y)$ in the domain,
\begin{gather}
\label{eq:shift}
D^n \sin(x,y) = \sin\Bigl(x + \frac{n\pi}{2},\; y - \frac{n\pi}{2}\Bigr), \\
\label{eq:shift-cos}
D^n \cos(x,y) = \cos\Bigl(x + \frac{n\pi}{2},\; y - \frac{n\pi}{2}\Bigr).
\end{gather}
In particular $D^2\sin(x,y) = -\sin(x,y)$ and $D^2\cos(x,y) = -\cos(x,y)$.
\end{thm}

\begin{proof}
Along the path, $\sin(x+t,y-t) = \sin(x+t)/\sin(x+y)$ with a constant
denominator, so $D^n\sin(x,y) = \sin(x+n\pi/2)/\sin(x+y)$, which is the
general sine at $(x + n\pi/2,\, y - n\pi/2)$ because the sum of the two
arguments is still $x+y$. The same argument applies to
$\cos(x,y) = \sin y/\sin(x+y)$, with $\frac{\mathrm d^n}{\mathrm dt^n}\sin(y-t)
= \sin(y - t - n\pi/2)$. For $n = 2$, Proposition~\ref{prop:symmetry} gives
$\sin(x+\pi, y-\pi) = -\sin(x,y)$ and $\cos(x+\pi,y-\pi) = -\cos(x,y)$.
\end{proof}

\begin{cor}[Primitives in a fixed frame]
\label{cor:primitives}
For every integer $n \ge 1$, the functions
$\sin(x - n\pi/2,\, y + n\pi/2)$ and $\cos(x - n\pi/2,\, y + n\pi/2)$
satisfy $D^n F = \sin(x,y)$ and $D^n F = \cos(x,y)$ respectively; every
$n$-th primitive along the lines $x+y = \text{const}$ differs from them by
a polynomial of degree less than $n$ in $t$ whose coefficients depend only
on $x+y$.
\end{cor}

\begin{proof}
Apply Theorem~\ref{thm:shift} to the shifted point. Two $n$-th primitives
along a line differ by a solution of $\mathrm{d}^n g/\mathrm{d}t^n = 0$.
\end{proof}

Theorem~\ref{thm:shift} says two things. With the frame held fixed, the
general cosine and sine obey the same equation $D^2 f = -f$ as the
classical pair. And one derivative is a quarter turn of the point:
$D\sin(x,y) = \sin(x+\pi/2,\,y-\pi/2)$ is the general sine of the point a
quarter turn further on the circle, read in the same frame. In the
notation \eqref{eq:CS}, this is the derivative of \eqref{eq:oblique-euler},
$\frac{\mathrm d}{\mathrm d\theta}\bigl(C_\theta + S_\theta u\bigr) = i\,e^{i\theta}
= C_{\theta+\pi/2} + S_{\theta+\pi/2}\,u$. The general tangent does not have this property, because both its
numerator and its denominator move along the path. Writing
$\tan(x+t,y-t) = \sin(x+y)\cot(y-t) - \cos(x+y)$ gives, for $n \ge 1$,
\[
D\tan(x,y) = \frac{\sin(x+y)}{\sin^2 y}, \qquad
D^n\tan(x,y) = \sin(x+y)\Bigl(-\frac{\mathrm d}{\mathrm dy}\Bigr)^{\!n-1}\!\csc^2 y .
\]

% =======================================================================
\section{Product identities and the General Sinus Transform}
\label{sec:products}
% =======================================================================

This section writes each general function as a finite product of copies
of itself at scaled arguments. The identities are exact and hold for any
number of factors. Replacing every factor by a ratio of two quadratic
polynomials then gives a rational approximation, which we call the
\emph{General Sinus Transform}.

The identities rest on the classical product of sines
\cite[\S 1.39]{GradshteynRyzhik2015}: for every integer $N \geq 1$ and
every real $\vartheta$,
\begin{equation}
\label{eq:product_of_sines}
\prod_{k=0}^{N-1} \sin\Bigl(\vartheta + \frac{\pi k}{N}\Bigr)
= \frac{\sin(N\vartheta)}{2^{N-1}} .
\end{equation}

\begin{thm}[Product identities]
\label{thm:products}
Let $N \geq 1$ be an integer and let $x, y \in \mathbb{R}$. If
$x + y \notin \pi\mathbb{Z}$, then
\begin{gather}
\label{eq:prod-sin}
\sin(x,y) = \prod_{k=0}^{N-1}\sin\Bigl(\frac{x + \pi k}{N},\, \frac{y}{N}\Bigr), \\
\label{eq:prod-cos}
\cos(x,y) = \prod_{k=0}^{N-1}\cos\Bigl(\frac{x}{N},\, \frac{y + \pi k}{N}\Bigr),
\end{gather}
and if $y \notin \pi\mathbb{Z}$, then
\begin{equation}
\label{eq:prod-tan}
\tan(x,y) = \prod_{k=0}^{N-1}\tan\Bigl(\frac{x + \pi k}{N},\, \frac{y + \pi k}{N}\Bigr).
\end{equation}
Every factor lies in the domain of its function.
\end{thm}

\begin{proof}
Apply \eqref{eq:product_of_sines} with $N\vartheta = x$, $N\vartheta = y$
and $N\vartheta = x+y$. The factors $2^{N-1}$ cancel in each quotient:
\begin{gather*}
\frac{\sin x}{\sin(x+y)} = \prod_{k=0}^{N-1}
\frac{\sin\bigl((x+\pi k)/N\bigr)}{\sin\bigl((x+\pi k)/N + y/N\bigr)}, \\
\frac{\sin y}{\sin(x+y)} = \prod_{k=0}^{N-1}
\frac{\sin\bigl((y+\pi k)/N\bigr)}{\sin\bigl(x/N + (y+\pi k)/N\bigr)}, \\
\frac{\sin x}{\sin y} = \prod_{k=0}^{N-1}
\frac{\sin\bigl((x+\pi k)/N\bigr)}{\sin\bigl((y+\pi k)/N\bigr)} .
\end{gather*}
By \eqref{eq:analytic} the left sides are $\sin(x,y)$, $\cos(x,y)$ and
$\tan(x,y)$, and each factor is the general function at the stated
arguments. No denominator vanishes, because the product of the
denominators is $\sin(x+y)/2^{N-1} \neq 0$, or $\sin y/2^{N-1} \neq 0$ for
the tangent.
\end{proof}

For $N = 2$ the first identity reads
$\sin(x,y) = \sin(\frac{x}{2},\frac{y}{2})\,\sin(\frac{x}{2}+\frac{\pi}{2},\frac{y}{2})$,
and repeating it $n$ times gives the case $N = 2^n$, the form in which the
identity appears in~\cite{GeneralSinus}.

As $N$ grows, the parts $x/N$, $y/N$ and $(x+y)/N$ of the arguments become
small while the shifts $\pi k/N$ stay exact. Write
$u_{N,k} = \sin(\pi k/N)$ and $v_{N,k} = \cos(\pi k/N)$. For small $s$,
$\sin(s + \pi k/N) = v_{N,k}\sin s + u_{N,k}\cos s \approx
v_{N,k}\,s + u_{N,k}(1 - s^2/2)$. With $s = z/N$ and after multiplication
by $-2N^2$, each sine in the factors of Theorem~\ref{thm:products} becomes
the quadratic
\begin{equation}
\label{eq:quadratic}
P_{N,k}(z) = u_{N,k}\,z^2 - 2N\,v_{N,k}\,z - 2N^2 u_{N,k} .
\end{equation}

\begin{defn}[General Sinus Transform]
\label{defn:GST}
The \emph{General Sinus Transform of order $N$} is the triple of rational
functions
\begin{gather*}
S_N(x,y) = \prod_{k=0}^{N-1}\frac{P_{N,k}(x)}{P_{N,k}(x+y)}, \qquad
C_N(x,y) = \prod_{k=0}^{N-1}\frac{P_{N,k}(y)}{P_{N,k}(x+y)}, \\
T_N(x,y) = \prod_{k=0}^{N-1}\frac{P_{N,k}(x)}{P_{N,k}(y)} .
\end{gather*}
\end{defn}

Each factor is a ratio of two quadratic polynomials, and $S_N$, $C_N$,
$T_N$ approximate $\sin(x,y)$, $\cos(x,y)$, $\tan(x,y)$ with products in
place of the sums of a truncated Taylor series. The approximation is not
exact for any finite $N$, since each factor carries the error of the
second-order expansion of $\sin s$ and $\cos s$. Surprisingly, the
approximation itself can be computed in closed form, because the
truncated factors are again sines, of slightly enlarged angles.

For $|\varepsilon| < \sqrt2$ let
\begin{equation}
\label{eq:rho-delta}
\rho(\varepsilon) = \sqrt{1 + \varepsilon^4/4}, \qquad
\delta(\varepsilon) = \arctan\frac{\varepsilon}{1 - \varepsilon^2/2},
\end{equation}
so that $1 - \varepsilon^2/2 = \rho\cos\delta$ and $\varepsilon = \rho\sin\delta$,
and $\delta(\varepsilon) = \varepsilon + \varepsilon^3/6 + O(\varepsilon^5)$.

\begin{thm}[Closed form of the transform]
\label{thm:GST-closed}
Let $x, y \in \mathbb{R}$ and let $N > \max(|x|, |y|, |x+y|)/\sqrt2$. Then,
whenever the denominators below do not vanish,
\begin{align}
\label{eq:SN-closed}
S_N(x,y) &= \Bigl(\frac{\rho(x/N)}{\rho((x+y)/N)}\Bigr)^{\!N}
\frac{\sin\bigl(N\delta(x/N)\bigr)}{\sin\bigl(N\delta((x+y)/N)\bigr)}, \\
C_N(x,y) &= \Bigl(\frac{\rho(y/N)}{\rho((x+y)/N)}\Bigr)^{\!N}
\frac{\sin\bigl(N\delta(y/N)\bigr)}{\sin\bigl(N\delta((x+y)/N)\bigr)}, \notag \\
T_N(x,y) &= \Bigl(\frac{\rho(x/N)}{\rho(y/N)}\Bigr)^{\!N}
\frac{\sin\bigl(N\delta(x/N)\bigr)}{\sin\bigl(N\delta(y/N)\bigr)}. \notag
\end{align}
\end{thm}

\begin{proof}
Let $\varepsilon = z/N$ and $\vartheta_k = \pi k/N$. By \eqref{eq:quadratic},
\begin{align*}
P_{N,k}(z) &= -2N^2\Bigl(u_{N,k}\bigl(1 - \tfrac{\varepsilon^2}{2}\bigr) + v_{N,k}\,\varepsilon\Bigr)
= -2N^2\rho(\varepsilon)\bigl(\sin\vartheta_k\cos\delta + \cos\vartheta_k\sin\delta\bigr) \\
&= -2N^2\rho(\varepsilon)\,\sin\bigl(\vartheta_k + \delta(\varepsilon)\bigr).
\end{align*}
By \eqref{eq:product_of_sines},
$\prod_{k=0}^{N-1} P_{N,k}(z) = (-2N^2\rho(\varepsilon))^N\sin(N\delta(\varepsilon))/2^{N-1}$.
Taking the quotients that define $S_N$, $C_N$ and $T_N$, the factors
$(-2N^2)^N/2^{N-1}$ cancel.
\end{proof}

The theorem says that the transform evaluates the general functions at
slightly perturbed angles: $S_N(x,y)$ is, up to the factor
$(\rho(x/N)/\rho((x+y)/N))^N = 1 + O(N^{-3})$, the general sine at
$\bigl(N\delta(x/N),\, N\delta((x+y)/N) - N\delta(x/N)\bigr)$, and
$N\delta(z/N) = z + z^3/(6N^2) + O(N^{-4})$. This gives the rate of
convergence exactly.

\begin{cor}[Rate of convergence]
\label{cor:GST-rate}
For fixed $x$, $y$ with $x + y \notin \pi\mathbb{Z}$, as $N \to \infty$,
\[
S_N(x,y) = \sin(x,y) + \frac{c(x,y)}{N^2} + O\bigl(N^{-3}\bigr),
\]
where
\[
c(x,y) = \frac{x^3\cos x\,\sin(x+y) - (x+y)^3\sin x\,\cos(x+y)}{6\sin^2(x+y)} .
\]
In particular $S_N \to \sin(x,y)$, and similarly $C_N \to \cos(x,y)$ and
$T_N \to \tan(x,y)$, each with an error of exact order $N^{-2}$ unless the
coefficient vanishes.
\end{cor}

\begin{proof}
Insert $N\delta(z/N) = z + z^3/(6N^2) + O(N^{-4})$ and
$\rho(z/N)^N = (1 + z^4/(4N^4))^{N/2} = 1 + O(N^{-3})$ into
\eqref{eq:SN-closed}, and expand $\sin(z + h) = \sin z + h\cos z + O(h^2)$ in
numerator and denominator.
\end{proof}

\begin{rem}[Numerical behaviour]
\label{rem:accuracy}
At $(x,y) = (0.3,2.7)$, close to the pole $x+y = \pi$, the coefficient is
$c = 66.14$, and the computed values of $N^2(S_N - \sin(x,y))$ are $66.28$
for $N = 64$ and $66.12$ for $N = 1000$. Over the four points
$(0.5,0.7)$, $(1,1)$, $(2,0.9)$ and $(0.3,2.7)$, the largest error of $S_N$
is about $1.6\times10^{-2}$ for $N = 64$, $6.3\times10^{-5}$ for $N = 1024$
and $2.5\times10^{-7}$ for $N = 16384$. The coefficient $c(x,y)$ grows like
$\sin^{-2}(x+y)$ near the poles, where $\sin(x,y)$ itself is unbounded. In
computations the values $u_{N,k}$, $v_{N,k}$ must be accurate: generating
them by a first-order rotation recurrence adds an error that accumulates
over the $N$ factors and dominates the error of the transform.
\end{rem}

The transform is of structural rather than computational interest:
evaluating $S_N$ costs $O(N)$ operations, far more than computing
$\sin x/\sin(x+y)$ directly. What it shows is that each general function
is, exactly, a product of copies of itself at small arguments
(Theorem~\ref{thm:products}), and that each copy is close to a ratio of
quadratic polynomials, much as the product formula
$\sin x = x\prod_{m\ge1}\bigl(1 - x^2/(m\pi)^2\bigr)$ expresses the
classical sine.

% =======================================================================
\section{Verification}
\label{sec:verification}
% =======================================================================

The script \texttt{verify\_gsp.py}, which accompanies the paper, checks
every numbered statement. Identities are evaluated at $50$ significant
digits at random points of their domains, with independent random draws for
each identity; derivatives, including $D^n$ for $n = 1, \dots, 6$ and the
primitives of Corollary~\ref{cor:primitives}, are computed symbolically with
SymPy and compared at $50$ digits; the product identities are tested for
every $N$ from $1$ to $40$; the closed form of
Theorem~\ref{thm:GST-closed} is compared with the products themselves and
the coefficient of Corollary~\ref{cor:GST-rate} with the computed errors;
the fixed-frame formulas of Proposition~\ref{prop:fixed-frame} are tested
for $n$ from $-3$ to $5$; and the numbers quoted in
Remark~\ref{rem:accuracy} are recomputed.

% =======================================================================
\section{Conclusion}
\label{sec:conclusion}
% =======================================================================

Projecting a point of the circle onto a line along an arbitrary direction,
instead of straight down, gives an oblique triangle, and the ratios of its
sides are the general sine, cosine and tangent. They are the coordinates
of the unit circle in an oblique frame
(Proposition~\ref{prop:coordinates}) and satisfy the oblique Pythagorean
identity \eqref{eq:pythagoras}. Their addition formulas, in the studied
and in the projection angle, all reduce to Ptolemy's identity
(Theorems~\ref{thm:add-studied} and~\ref{thm:add-projection}).
With the frame held fixed, they obey addition formulas with a single
cross term and De Moivre's formula (Proposition~\ref{prop:fixed-frame}),
and differentiation turns the point by a quarter turn, so that
$D^2 f = -f$ (Theorem~\ref{thm:shift}). They are exact products of their
own copies at scaled arguments, for any number of factors
(Theorem~\ref{thm:products}); truncating each factor gives a rational
approximation which is itself a ratio of sines in closed form, and whose
error is of exact order $N^{-2}$ (Theorem~\ref{thm:GST-closed} and
Corollary~\ref{cor:GST-rate}).

Two questions remain open. The first is whether the general functions
satisfy orthogonality relations, for a suitable weight, that parallel
those of the classical functions; with the frame fixed, $C_\theta$ and
$S_\theta$ are combinations of $\cos\theta$ and $\sin\theta$, but the natural
inner product of the oblique frame is not the orthonormal one. The second
is the analogue of the construction for the hyperbolic functions, where
the oblique frame of the Euclidean plane is replaced by a frame of the
split-complex plane~\cite{Yaglom1968,HarkinHarkin2004}.

\subsection*{Declarations}

{\sloppy
\noindent\textbf{Funding.} No funding was received for this work.

\noindent\textbf{Competing interests.} The author declares no competing interests.

\noindent\textbf{Data and code availability.} No data were used. The
verification script, \texttt{verify\_gsp.py}, is provided as
supplementary material with the submission.

\noindent\textbf{Declaration of generative AI use.} Claude (Anthropic), including
Claude Opus 5.5, helped draft and edit the text and check calculations;
the mathematics is the author's own.\par}

\begin{thebibliography}{99}

\bibitem{CoxeterGreitzer1967}
Coxeter, H.S.M., Greitzer, S.L.: \textit{Geometry Revisited}.
New Mathematical Library, vol.~19. Mathematical Association of America,
Washington, D.C. (1967)

\bibitem{GeneralSinus}
Derraz, M.: General Sinus.
\textit{Int. J. Swarm Intell. Evol. Comput.} \textbf{6}(2) (2017).
\url{https://doi.org/10.4172/2090-4908.1000155}

\bibitem{DerrazGN}
Derraz, M.: General Numbers: the number plane beyond the right angle.
Zenodo preprint (2026). \url{https://doi.org/10.5281/zenodo.24091305}

\bibitem{GradshteynRyzhik2015}
Gradshteyn, I.S., Ryzhik, I.M.:
\textit{Table of Integrals, Series, and Products}, 8th edn.
Academic Press, Amsterdam (2015)

\bibitem{HarkinHarkin2004}
Harkin, A.A., Harkin, J.B.: Geometry of generalized complex numbers.
\textit{Math. Mag.} \textbf{77}(2), 118--129 (2004)

\bibitem{KantorSolodovnikov1989}
Kantor, I.L., Solodovnikov, A.S.:
\textit{Hypercomplex Numbers: An Elementary Introduction to Algebras}.
Springer, New York (1989)

\bibitem{Lindqvist1995}
Lindqvist, P.: Some remarkable sine and cosine functions.
\textit{Ricerche Mat.} \textbf{44}(2), 269--290 (1995)

\bibitem{Maor1998}
Maor, E.: \textit{Trigonometric Delights}.
Princeton University Press, Princeton (1998)

\bibitem{Yaglom1968}
Yaglom, I.M.: \textit{Complex Numbers in Geometry}.
Academic Press, New York (1968)

\end{thebibliography}

\end{document}
